\documentclass[../../main.tex]{subfiles}% 注意这里的文件路径不能用 ./main.tex ,否则用latexmk编译子文件会报错
\graphicspath{{\subfix{./image/}}{\subfix{../../image/}}} % 指定图片引用路径,后续可以直接使用图片文件名
% 如果图片存放在上级目录,则使用 ../../image/ 作为备用路径

% 例如:
% \begin{figure}[H]
% \centering
% \includegraphics[width=0.3\textwidth]{图.png}
% \caption{}
% \label{figure:图}
% \end{figure}
% 注意:上述\label{}一定要放在\caption{}之后,否则引用图片序号会只会显示??.

\begin{document}

\section{场论}

\begin{definition}
设$D\subseteq\mathbb{R}^n$为非空开集，称映射
\begin{align*}
f:D\to\mathbb{R},\,\,\bm{x}=(x_1,\cdots,x_n)\mapsto f(\bm{x})
\end{align*}
为$D$上的\textbf{标量场}。
称映射
\begin{align*}
F:D\to\mathbb{R}^m,\,\,\bm{x}=(x_1,\cdots,x_n)\mapsto \bm{F}(\bm{x})
\end{align*}
为$D$上的\textbf{向量场}。通常取$m=n$，即
\begin{align*}
F:D\to\mathbb{R}^n,\,\,\bm{x}=(x_1,\cdots,x_n)\mapsto \bm{F}(\bm{x}).
\end{align*}
\end{definition}

\begin{definition}
设$f:D\subseteq\mathbb{R}^n\to\mathbb{R}$在点$\bm{x}_0$处所有一阶偏导数存在，称向量
\[
\nabla f(\bm{x}_0)=\left(\frac{\partial f}{\partial x_1}(\bm{x}_0),\cdots,\frac{\partial f}{\partial x_n}(\bm{x}_0)\right)
\]
为$f$在$\bm{x}_0$处的\textbf{梯度},称$\nabla$为$\mathbf{Hamilton}$\textbf{(哈密顿)算子}.
\end{definition}

\begin{proposition}\label{proposition:连通开集上梯度为0就是常数}
设 $\Omega\subset \mathbb{R}^n$ 是非空连通开集，$f:\Omega\to\mathbb{R}$ 在 $\Omega$ 上可微，且
\begin{align*}
\nabla f(x)=0,\qquad \forall x\in\Omega,
\end{align*}
则 $f$ 在 $\Omega$ 上为常数.
\end{proposition}
\begin{remark}
该命题的逆命题显然成立：常函数的梯度处处为零.
\end{remark}
\begin{proof}

任取 $x_0\in\Omega$，定义集合
\begin{align*}
A\triangleq\{\,x\in\Omega\mid f(x)=f(x_0)\,\}.
\end{align*}
显然$x_0\in A$，故$A\neq\varnothing$.

先证 $A$ 是开集：任取 $y\in A$，因 $\Omega$ 开，故存在 $r>0$ 使得 $B(y,r)\subset\Omega$. 对任意 $z\in B(y,r)$，由于球是凸的，线段 $[y,z]\subset B(y,r)\subset\Omega$. 对 $t\in[0,1]$，令 $\gamma(t)=y+t(z-y)$，则由链式法则得
\begin{align*}
\frac{\mathrm{d}}{\mathrm{d}t} f(\gamma(t))=\nabla f(\gamma(t))\cdot(z-y)=0.
\end{align*}
故 $f(\gamma(t))$ 恒为常数，因此$f(\gamma(0))=f(\gamma(1))=f(z)=f(y)=f(x_0)$，所以 $z\in A$. 因此 $B(y,r)\subset A$，即$A$ 为开集.

再证 $A$ 是闭集：由 $f$ 的连续性知$A=f^{-1}(\{f(x_0)\})$ 是 $\Omega$ 的闭集.

若$A \subsetneq  \Omega$,则$\Omega =A\sqcup (\Omega\setminus A)$,由$A$ 非空且既开又闭知$\Omega\setminus A$也是开集,这与$\Omega$ 连通矛盾!故 $A=\Omega$. 因此对任意 $x\in\Omega$，$f(x)=f(x_0)$，即 $f$ 为常数.
\end{proof}

\begin{proposition}\label{proposition:}
设 $A$ 是 $n$ 阶实矩阵，$x=(x_1,\cdots,x_n)^T\in\mathbb{R}^n$，定义二次型
\begin{align*}
f(x)\triangleq x^T A x=\sum_{i=1}^n\sum_{j=1}^n a_{ij}x_i x_j.
\end{align*}
则 $f$ 在 $x$ 处的梯度为
\begin{align}
\nabla f(x)=(A+A^T)x. \label{grad:general}
\end{align}
特别地，若 $A$ 为对称矩阵，则
\begin{align}
\nabla f(x)=2Ax. \label{grad:symmetric}
\end{align}
\end{proposition}
\begin{remark}
二次型 $x^T A x$ 的值仅由 $A$ 的对称部分 $\frac{1}{2}(A+A^T)$ 决定，因此梯度的表达式自然地反映了这一事实.
\end{remark}
\begin{proof}
将 $f(x)$ 展开：
\begin{align*}
f(x)=\sum_{i=1}^n\sum_{j=1}^n a_{ij}x_i x_j.
\end{align*}
固定 $k$，对 $x_k$ 求偏导：
\begin{align*}
\frac{\partial f}{\partial x_k}
=\sum_{i=1}^n\sum_{j=1}^n a_{ij}\frac{\partial}{\partial x_k}(x_i x_j) 
=\sum_{i=1}^n\sum_{j=1}^n a_{ij}(\delta_{ik}x_j+x_i\delta_{jk}) 
=\sum_{j=1}^n a_{kj}x_j+\sum_{i=1}^n a_{ik}x_i.
\end{align*}
上式第一项为 $(Ax)_k$，第二项为 $(A^T x)_k$，因此
\begin{align*}
\frac{\partial f}{\partial x_k}=((A+A^T)x)_k.
\end{align*}
所以 $\nabla f(x)=(A+A^T)x$，即 \eqref{grad:general}式成立.

当 $A$ 对称时，$A^T=A$，代入 \eqref{grad:general} 得 $\nabla f(x)=2Ax$，即 \eqref{grad:symmetric}式成立.
\end{proof}

\begin{definition}
设 $A=(a_{ij})_{m\times n}$ 为复矩阵，其中 $a_{ij}\in\mathbb C$ 为独立复变量。固定指标 $j\in\{1,\cdots,m\}$，$k\in\{1,\cdots,n\}$，将矩阵 $A$ 视为仅关于标量变量 $a_{jk}$ 的函数（其余元素视为常数）。则 $A$ 对 $a_{jk}$ 的偏导数定义为逐元素求导得到的矩阵
\begin{align*}
\frac{\partial A}{\partial a_{jk}}
\triangleq
\begin{pmatrix}
\dfrac{\partial a_{11}}{\partial a_{jk}} & \dfrac{\partial a_{12}}{\partial a_{jk}} & \cdots & \dfrac{\partial a_{1n}}{\partial a_{jk}} \\[1.2ex]
\dfrac{\partial a_{21}}{\partial a_{jk}} & \dfrac{\partial a_{22}}{\partial a_{jk}} & \cdots & \dfrac{\partial a_{2n}}{\partial a_{jk}} \\
\vdots & \vdots & \ddots & \vdots \\
\dfrac{\partial a_{m1}}{\partial a_{jk}} & \dfrac{\partial a_{m2}}{\partial a_{jk}} & \cdots & \dfrac{\partial a_{mn}}{\partial a_{jk}}
\end{pmatrix}.
\end{align*}
\end{definition}
\begin{remark}
上述定义中，若矩阵元素为复数，则偏导数按实部和虚部分别求导处理，所得结果仍为复矩阵。
\end{remark}

\begin{theorem}[矩阵偏导运算的基本性质]
\label{theorem:矩阵偏导运算的基本性质}
\begin{enumerate}[(1)]
\item 设 $I\subseteq\mathbb R$ 为开区间，$t\in I$ 为实自变量。设
\begin{align*}
A(t)=(a_{ij}(t))_{m\times n},\qquad B(t)=(b_{ij}(t))_{n\times p},
\end{align*}
其中每个元素 $a_{ij}:I\to\mathbb C$，$b_{ij}:I\to\mathbb C$ 均为关于 $t$ 的可微复值函数。则有
\[
\frac{\mathrm{d}}{\mathrm{d}t}\big(A(t)B(t)\big)=\frac{\mathrm{d}A\left( t \right)}{\mathrm{d}t}B(t)+A(t)\frac{\mathrm{d}B\left( t \right)}{\mathrm{d}t}.
\]
特别地，若$x=(x_1(t),\cdots,x_n(t))^T$，则有
\[
\frac{\mathrm{d}}{\mathrm{d}t}\big(A(t)x(t)\big)=\frac{\mathrm{d}A\left( t \right)}{\mathrm{d}t}x(t)+A(t)\frac{dx(t)}{dt}.
\]
进而若函数 \(f: \mathbb{R}^n \to \mathbb{R}\) 在点 \(\boldsymbol{x}_0\) 的某邻域内二阶连续可微，给定方向向量 \(\boldsymbol{h} \in \mathbb{R}^n\)。定义一元函数 \(\varphi(t) = f(\boldsymbol{x}_0 + t\boldsymbol{h})\)，则根据多元复合函数求导的链式法则，\(\varphi(t)\) 对 \(t\) 的导数满足
\begin{equation*}
\varphi'(t) = \frac{\mathrm{d}}{\mathrm{d}t} \left[ \boldsymbol{h} \cdot \nabla f(\boldsymbol{x}_0 + t\boldsymbol{h}) \right] = \boldsymbol{h}^T H(\boldsymbol{x}_0 + t\boldsymbol{h}) \boldsymbol{h}
\end{equation*}
其中\(H(\boldsymbol{x})\) 为 \(f\) 在点 \(\boldsymbol{x}\) 处的Hess矩阵.

\item 设$A=(a_{ij})_{n\times m}$是复数域上的矩阵,对于矩阵 \(A = (a_{ij})\)的每个元素，关于其元素 \(a_{jk}\) 求偏导，有
\[
\frac{\partial A}{\partial a_{jk}} = E_{jk}.
\]
\end{enumerate}
\end{theorem}
\begin{proof}
\begin{enumerate}[(1)]
\item 考察矩阵乘积 \(AB\) 的第 \(i\) 行第 \(q\) 列元素：
\[
(AB)_{iq}=\sum_{\nu=1}^{n} a_{i\nu}(t)b_{\nu q}(t).
\]
对 \(t\) 求导得
\[
\frac{\mathrm{d}}{\mathrm{d}t}(AB)_{iq}=\sum_{\nu =1}^n{\left( \frac{\mathrm{d}a_{i\nu}}{\mathrm{d}t}b_{\nu q}+a_{i\nu}\frac{\mathrm{d}b_{\nu q}}{\mathrm{d}t} \right)}.
\]
右端第一项为 \(\left(\frac{\mathrm{d}A}{\mathrm{d}t}B\right)_{iq}\)，第二项为 \(\left(A\frac{\mathrm{d}B}{\mathrm{d}t}\right)_{iq}\)。由于 \(i,q\) 任意，故矩阵等式成立。

特别地，只需取$l=1$,此时$B$就是$n$维列向量,由上述结论立得.

\item 对矩阵 \(A\) 的第 \(p\) 行第 \(q\) 列元素 \(a_{pq}\)关于 \(a_{jk}\) 求偏导，得
\[
\frac{\partial a_{pq}}{\partial a_{jk}} = \delta_{pj} \delta_{qk}.
\]
其中$\delta$是Kronecker符号.
因此，导数矩阵的第 \((p,q)\) 个分量仅在 \(p=j\) 且 \(q=k\) 时为 \(1\)，其余为 \(0\)。这正是标准基矩阵 \(E_{jk}\) 的定义。故
\[
\frac{\partial A}{\partial a_{jk}} = E_{jk}.
\]
\end{enumerate}
\end{proof}

\begin{example}
设 \(A = (a_{ij}) \in \mathbb{R}^{n \times n}\) 为可逆矩阵，\(b \in \mathbb{R}^{n}\) 为常向量。若线性方程组
\[
Ax = b
\]
的解向量为\(x = (x_1, \cdots, x_n)^T\),且记 \(A^{-1} = (c_{ij})\)。则对于任意的 \(i, j, k \in \{1, \cdots, n\}\)，有
\[
\frac{\partial x_i}{\partial a_{jk}} = -c_{ij} x_k.
\]
\end{example}
\begin{proof}
由于 \(b\) 为常数，对方程 \(Ax = b\) 两边关于 \(a_{jk}\) 求偏导。左边应用\cref{theorem:矩阵偏导运算的基本性质}，右边导数为 \(0\)，可得：
\[
E_{jk} x + A \frac{\partial x}{\partial a_{jk}} = 0\Longrightarrow A \frac{\partial x}{\partial a_{jk}} = -E_{jk} x\Longrightarrow \frac{\partial x}{\partial a_{jk}} = -A^{-1} E_{jk} x.
\]
注意 \(E_{jk} x\) 是一个列向量，其第 \(j\) 个分量为 \(x_k\)，其余分量为 \(0\)。因此乘积 \(A^{-1} (E_{jk} x)\) 就是取 \(A^{-1}\) 的第 \(j\) 列乘以标量 \(x_k\)。而 \(A^{-1}\) 的第 \(j\) 列元素恰好为 \(c_{1j}, c_{2j}, \cdots, c_{nj}\)。
所以，该乘积向量的第 \(i\) 个分量为：
\[
\left( A^{-1} E_{jk} x \right)_i = c_{ij} x_k.
\]
代入上式，取第 \(i\) 个分量，得到
\[
\frac{\partial x_i}{\partial a_{jk}} = -c_{ij} x_k.
\]\end{proof}

\begin{definition}
设$\bm{F}=(F_1,\cdots,F_n)$是$\mathbb{R}^n$上的向量场，若$F_i$均可微，则定义$\bm{F}$的\textbf{散度}为
\[
\nabla\cdot\bm{F}=\sum_{i=1}^{n}\frac{\partial F_i}{\partial x_i}.
\]
\end{definition}

\begin{definition}
设$\bm{F}=(F_1,F_2,F_3)$是$\mathbb{R}^3$上的向量场，若各分量均可微，则定义$\bm{F}$的\textbf{旋度}为
\[
\mathrm{rot}\,\boldsymbol{F}=\nabla\times\bm{F}=
\begin{vmatrix}
\bm{i} & \bm{j} & \bm{k} \\
\frac{\partial}{\partial x} & \frac{\partial}{\partial y} & \frac{\partial}{\partial z} \\
F_1 & F_2 & F_3
\end{vmatrix}
=\left(\frac{\partial F_3}{\partial y}-\frac{\partial F_2}{\partial z},\frac{\partial F_1}{\partial z}-\frac{\partial F_3}{\partial x},\frac{\partial F_2}{\partial x}-\frac{\partial F_1}{\partial y}\right).
\]
\end{definition}

\begin{definition}
对于二次可微的标量场$f:\mathbb{R}^n\to\mathbb{R}$，定义\textbf{Laplace(拉普拉斯)算子}为
\[
\Delta f=\nabla\cdot(\nabla f)=\sum_{i=1}^{n}\frac{\partial^{2}f}{\partial x_i^{2}}.
\]
\end{definition}

\begin{definition}
设$D\subseteq\mathbb{R}^n$为开集，$f:D\to\mathbb{R}$，$\bm{x}_0\in D$，$\bm{e}\in\mathbb{R}^n$为单位向量。若极限
\begin{align*}
\lim_{t\to 0}\frac{f(\bm{x}_0+t\bm{e})-f(\bm{x}_0)}{t}
\end{align*}
存在，则称该极限为$f$在点$\bm{x}_0$处沿方向$\bm{e}$的\textbf{方向导数}，记作$\frac{\partial f}{\partial\bm{e}}(\bm{x}_0)$或$D_{\bm{e}}f(\bm{x}_0)$.
\end{definition}

\begin{theorem}[方向导数的基本性质]\label{theorem:方向导数的基本性质}
设$f,g:D\subseteq\mathbb{R}^n\to\mathbb{R}$在点$\bm{x}_0\in D$处可微，$\alpha,\beta\in\mathbb{R}$，$\bm{e}\in \mathbb{R}^n$为任意单位向量，$\nabla f(\bm{x}_0)$为梯度。则
\begin{enumerate}[(1)]
\item\label{theorem:方向导数的基本性质-1} \textbf{梯度表示:}方向导数存在且
\begin{align*}
\frac{\partial f}{\partial\bm{e}}(\bm{x}_0)=\nabla f(\bm{x}_0)\cdot\bm{e}.
\end{align*}

\item \textbf{线性性质:}
\begin{align*}
\frac{\partial(\alpha f+\beta g)}{\partial\bm{e}}(\bm{x}_0)=\alpha\frac{\partial f}{\partial\bm{e}}(\bm{x}_0)+\beta\frac{\partial g}{\partial\bm{e}}(\bm{x}_0).
\end{align*}

\item 若$\nabla f(\bm{x}_0)\neq\bm{0}$，则当$\bm{e}$与$\nabla f(\bm{x}_0)$同向时方向导数取得最大值$|\nabla f(\bm{x}_0)|$，反向时取得最小值$-|\nabla f(\bm{x}_0)|$.进而对所有单位方向$\bm{e}$有
\begin{align*}
-|\nabla f(\bm{x}_0)|\leqslant\frac{\partial f}{\partial\bm{e}}(\bm{x}_0)\leqslant|\nabla f(\bm{x}_0)|.
\end{align*}

\item 设$\varphi(t)=f(\bm{x}_0+t\bm{e})$，则$\varphi'(0)=\frac{\partial f}{\partial\bm{e}}(\bm{x}_0)$.并且当$f$可微时,有$\varphi'(t)=\nabla f(\bm{x}_0+t\bm{e})\cdot\bm{e}$.
\end{enumerate}
\end{theorem}
\begin{remark}
方向导数存在不能推出函数在该点连续，即使所有方向导数存在且为零，函数仍可能不连续。例如
\[
f(x,y)=\begin{cases}\dfrac{x^{2}y}{x^{4}+y^{2}},&(x,y)\neq(0,0),\\0,&(x,y)=(0,0),\end{cases}
\]
在原点沿任何方向的方向导数均为$0$，但$f$在原点不连续.
\end{remark}

\begin{theorem}\label{theorem:散度旋度Laplace算子的性质}
设$f,g$是标量场，$\bm{F},\bm{G}$是向量场，$c\in \mathbb{R}$，则
\begin{enumerate}[(1)]
\item\label{theorem:散度旋度Laplace算子的性质-1} $\nabla(cf)=c\nabla f$，$\quad \nabla\cdot(c\bm{F})=c\nabla\cdot\bm{F}$，$\quad \nabla\times(c\bm{F})=c\nabla\times\bm{F}$.

\item\label{theorem:散度旋度Laplace算子的性质-2} $\nabla(f\pm g)=\nabla f\pm\nabla g$，$\quad \nabla\cdot(\bm{F}\pm\bm{G})=\nabla\cdot\bm{F}\pm\nabla\cdot\bm{G}$，$\quad \nabla\times(\bm{F}\pm\bm{G})=\nabla\times\bm{F}\pm\nabla\times\bm{G}$.

\item\label{theorem:散度旋度Laplace算子的性质-3} $\nabla(fg)=f\nabla g+g\nabla f$.

\item\label{theorem:散度旋度Laplace算子的性质-4} $\nabla\cdot(f\bm{F})=f\nabla\cdot\bm{F}+\nabla f\cdot\bm{F}$.

\item\label{theorem:散度旋度Laplace算子的性质-5} $\nabla\times(f\bm{F})=f\nabla\times\bm{F}+\nabla f\times\bm{F}$.

\item\label{theorem:散度旋度Laplace算子的性质-6} $\nabla\cdot(\bm{F}\times\bm{G})=\bm{G}\cdot(\nabla\times\bm{F})-\bm{F}\cdot(\nabla\times\bm{G})$.

\item\label{theorem:散度旋度Laplace算子的性质-7} 旋度无源:$\nabla\times(\nabla f)=\bm{0}$.

\item\label{theorem:散度旋度Laplace算子的性质-8} 散度无旋:$\nabla\cdot(\nabla\times\bm{F})=0$.

\item\label{theorem:散度旋度Laplace算子的性质-9} $\Delta(cu)=c\Delta u$
\item $\Delta(u\pm v)=\Delta u\pm\Delta v$.

\item\label{theorem:散度旋度Laplace算子的性质-10} $\Delta(uv)=u\Delta v+v\Delta u+2\nabla u\cdot\nabla v$,特别地$\Delta u^2=2u\Delta u+\left| \nabla u \right|^2.$
\end{enumerate}
\end{theorem}

\begin{theorem}[Gauss-Green公式(散度定理)]\label{theorem:Gauss-Green公式(散度定理)}
设$\Omega$ 是 $\mathbb{R}^n$ 中的一个有界开集，且 $\partial\Omega\in C^1$。若向量场 $\bm{F}=(F_1,F_2,\cdots,F_n):\Omega\to\mathbb{R}^n$且$\bm{F}\in C^1(\Omega)\cap C(\overline{\Omega})$，则
\begin{align*}
\int_\Omega \operatorname{div} \bm{F}\,\mathrm{d}x =\int_{\Omega}\nabla\cdot\bm{F}\,\mathrm{d}V= \int_{\partial\Omega} \bm{F}\cdot \bm{n}\,\mathrm{d}S(x).
\end{align*}
其中 $\bm{n}$ 为 $\partial\Omega$ 的单位外法向量。
\end{theorem}

\begin{theorem}[Green公式（平面散度定理与旋度定理）]\label{theorem:Green公式平面散度定理与旋度定理}
设$D\subset\mathbb{R}^2$是有界闭区域，边界$\partial D$是分段光滑的闭曲线，取正向（逆时针）。若$P,Q$在$\overline{D}$上连续，在$D$内连续可微，则
\begin{align*}
\iint_{D}\left(\frac{\partial Q}{\partial x}-\frac{\partial P}{\partial y}\right)\mathrm{d}x\mathrm{d}y=\oint_{\partial D}P\,\mathrm{d}x+Q\,\mathrm{d}y, 
\end{align*}
\begin{align}
\iint_{D}\left(\frac{\partial P}{\partial x}+\frac{\partial Q}{\partial y}\right)\mathrm{d}x\mathrm{d}y=\oint_{\partial D}(P\,\mathrm{d}y-Q\,\mathrm{d}x).\label{eq:green_div} 
\end{align}
\end{theorem}
\begin{remark}
上述定理中的\eqref{eq:green_div}是散度定理在二维的体现.
\end{remark}

\begin{theorem}[Green恒等式]\label{theorem:Green恒等式}
设$\Omega\subset\mathbb{R}^n$是有界闭区域，其边界$\partial\Omega$是逐片光滑的定向曲面，$\bm{n}$为$\partial\Omega$的单位外法向量。
\begin{enumerate}[(1)]
\item\label{theorem:Green第一恒等式} $\mathbf{Green}$\textbf{第一恒等式:}若$u,v\in C^2(\Omega)\cap C^1(\overline{\Omega})$，则有
\begin{align*}
\int_{\Omega}u\Delta v\,\mathrm{d}V=\int_{\partial\Omega}u\frac{\partial v}{\partial\bm{n}}\,\mathrm{d}S-\int_{\Omega}\nabla u\cdot\nabla v\,\mathrm{d}V.
\end{align*}
特别地,取$u\equiv 1$,则有
\begin{align*}
\int_{\Omega}\Delta v\,\mathrm{d}V=\int_{\partial\Omega}\frac{\partial v}{\partial\bm{n}}\,\mathrm{d}S.
\end{align*}

\item\label{theorem:Green第二恒等式} $\mathbf{Green}$\textbf{第二恒等式:}若$u,v\in C^2(\Omega)\cap C^1(\overline{\Omega})$，则有
\begin{align*}
\int_{\Omega}(u\Delta v-v\Delta u)\,\mathrm{d}V=\int_{\partial\Omega}\left(u\frac{\partial v}{\partial\bm{n}}-v\frac{\partial u}{\partial\bm{n}}\right)\mathrm{d}S. 
\end{align*}

\item\label{theorem:Green恒等式-3} 若$u,v\in C^1(\Omega)\cap C(\overline{\Omega})$，则有
\begin{align*}
\int_\Omega \frac{\partial u}{\partial x_i}v\,\mathrm{d}x = -\int_\Omega u \frac{\partial v}{\partial x_i}\,\mathrm{d}x + \int_{\partial\Omega} u v n_i\,\mathrm{d}S(x),
\end{align*}
其中 $n_i$ 是 $\partial\Omega$ 的单位外法向量 $\bm{n}$ 的第 $i$ 个分量.特别地，取$v\equiv 1$，则有
\begin{align*}
\int_{\Omega}\frac{\partial u}{\partial x_i}\,\mathrm{d}V=\int_{\partial\Omega}u\,n_i\,\mathrm{d}S. 
\end{align*}
\end{enumerate}
\end{theorem}
\begin{proof}
\begin{enumerate}[(1)]
\item 取 $\bm{F} = u \nabla v$，其中 $u, v \in C^2(\Omega) \cap C^1(\overline{\Omega})$。则由\rrefthe{theorem:散度旋度Laplace算子的性质}{theorem:散度旋度Laplace算子的性质-4}和Laplace算子的定义知
\[
\nabla \cdot \bm{F} = \nabla \cdot (u \nabla v) = \nabla u \cdot \nabla v + u \nabla \cdot (\nabla v) = \nabla u \cdot \nabla v + u \Delta v.
\]
又由$u, \nabla v \in C^1(\overline{\Omega})$知$\bm{F} \in C^1(\overline{\Omega};\mathbb{R}^n)$，故由\hyperref[theorem:Gauss-Green公式(散度定理)]{散度定理}和\hyperref[theorem:方向导数的基本性质-1]{方向导数的基本性质}得
\[
\int_{\Omega}\nabla\cdot\bm{F}\,\mathrm{d}V=\int_{\partial\Omega}\bm{F}\cdot\bm{n}\,\mathrm{d}S \iff \int_{\Omega} (u \Delta v + \nabla u \cdot \nabla v) \, \mathrm{d}V = \int_{\partial\Omega} (u \nabla v) \cdot \bm{n} \, \mathrm{d}S = \int_{\partial\Omega} u \frac{\partial v}{\partial \bm{n}} \, \mathrm{d}S,
\]
其中 $\frac{\partial v}{\partial \bm{n}} = \nabla v \cdot \bm{n}$ 是 $v$ 沿外法向$\bm{n}$的方向导数。移项即得
\begin{align}\label{eq::HJIOUJWFIOJO3424288}
\int_{\Omega} u \Delta v \, \mathrm{d}V = \int_{\partial\Omega} u \frac{\partial v}{\partial \bm{n}} \, \mathrm{d}S - \int_{\Omega} \nabla u \cdot \nabla v \, \mathrm{d}V.
\end{align}

\item 将\eqref{eq::HJIOUJWFIOJO3424288}式中的 $u$ 与 $v$ 互换，得到
\[
\int_{\Omega} v \Delta u \, \mathrm{d}V = \int_{\partial\Omega} v \frac{\partial u}{\partial \bm{n}} \, \mathrm{d}S - \int_{\Omega} \nabla v \cdot \nabla u \, \mathrm{d}V.
\]
将\eqref{eq::HJIOUJWFIOJO3424288}式减去上式，再结合$\nabla u \cdot \nabla v = \nabla v \cdot \nabla u$，便有
\[
\int_{\Omega} (u \Delta v - v \Delta u) \, \mathrm{d}V = \int_{\partial\Omega} \left( u \frac{\partial v}{\partial \bm{n}} - v \frac{\partial u}{\partial \bm{n}} \right) \mathrm{d}S.
\]

\item 取向量场
\begin{align*}
\bm{F}=\left( 0,\cdots, 0, uv, 0, \cdots, 0 \right),
\end{align*}
其中非零分量在第$i$个位置, 即$\bm{F}=uv\bm{e}_i$. 由\hyperref[theorem:Gauss-Green公式(散度定理)]{散度定理}得
\begin{align*}
\int_{\Omega}{\nabla \cdot \mathbf{F}\,\mathrm{d}x}=\int_{\partial \Omega}{\mathbf{F}\cdot \mathbf{n}\,\mathrm{d}S}&\Longleftrightarrow \int_{\Omega}{\frac{\partial}{\partial x_i}\left( uv \right) \mathrm{d}x}=\int_{\Omega}{\left( \frac{\partial u}{\partial x_i}v+u\frac{\partial v}{\partial x_i} \right) \mathrm{d}x}=\int_{\partial \Omega}{uvn_i\,\mathrm{d}S}
\\
&\Longleftrightarrow \int_{\Omega}{\frac{\partial u}{\partial x_i}v\mathrm{d}x}=-\int_{\Omega}{u\frac{\partial v}{\partial x_i}\mathrm{d}x}+\int_{\partial \Omega}{uvn_i\,\mathrm{d}S}.
\end{align*}
\end{enumerate}
\end{proof}

\begin{theorem}[多元函数分部积分公式]\label{theorem:多元函数分部积分公式}
设$\Omega\subset\mathbb{R}^n$是有界闭区域，其边界$\partial\Omega$是逐片光滑的定向曲面，$\bm{n}$为$\partial\Omega$的单位外法向量。若$u\in C^1(\Omega)\cap C(\overline{\Omega})$，$\bm{F}\in C^1(\Omega)\cap C(\overline{\Omega})$为向量场，则
\begin{align*}
\int_{\Omega}u(\nabla\cdot\bm{F})\,\mathrm{d}V=\int_{\partial\Omega}u(\bm{F}\cdot\bm{n})\,\mathrm{d}S-\int_{\Omega}\nabla u\cdot\bm{F}\,\mathrm{d}V. 
\end{align*}
\end{theorem}

\begin{definition}[径向对称函数]
设$f:\mathbb{R}^n\to \mathbb{R}$,若存在一个函数$g:[0,+\infty)\to \mathbb{R}$,使得对$\forall x\in \mathbb{R}^n$,都有
\begin{align*}
f(x)=g(|x|),
\end{align*}
其中$|x|$是欧氏空间$\mathbb{R}^n$的欧氏范数.此时我们称$f$是\textbf{径向对称的}.
\end{definition}

\begin{example}
设 $f(x)$ 为定义在 $\mathbb{R}^3$ 上的连续函数，且满足
\begin{enumerate}[(1)]
\item $f(x)$ 径向对称;
\item $\lim\limits_{|x|\to+\infty}f(x)=0$;
\item $f(x)$ 在 $\mathbb{R}^3$ 上具有连续偏导数，且
\[
\iiint_{\mathbb{R}^3}|\nabla f(x)|\,\mathrm{d}x<\infty.
\]
\end{enumerate}
证明：存在不依赖于 $f(x)$ 的正常数 $C$，使得
\[
\max_{x\in\mathbb{R}^3}\bigl(|x|^2|f(x)|\bigr)
\leqslant C\iiint_{\mathbb{R}^3}|\nabla f(x)|\,\mathrm{d}x.
\]
\end{example}
\begin{proof}
由$f$ 径向对称知,存在一个函数$g:[0,+\infty)\to\mathbb{R}$,使得
\[
f(x)=g(|x|)\quad \forall x\in\mathbb{R}^3.
\]
当$x\in\mathbb{R}^3\setminus \{0\}$时，记$r=|x|>0$. 由微积分基本定理，对$\forall R>r$，有
\[
g(R)-g(r)=\int_r^R g'(s)\,\mathrm{d}s.
\]
令 $R\to+\infty$，并利用 $\lim\limits_{|x|\to+\infty}f(x)=0$得
\[
\lim\limits_{R\to+\infty}g(R)=\lim\limits_{R\to+\infty}f((R,0,0))=0 \Longrightarrow g(r)=-\int_r^{+\infty}g'(s)\,\mathrm{d}s.
\]
因此
\begin{align}
|g(r)|\leqslant \int_r^{+\infty}|g'(s)|\,\mathrm{d}s\Longrightarrow r^2|g(r)|\leqslant r^2\int_r^{+\infty}|g'(s)|\,\mathrm{d}s \leqslant \int_r^{+\infty}|g'(s)|s^2\,\mathrm{d}s.\label{eq::iojifjioJIJIJw-1}
\end{align}
记 $r=|x|=\sqrt{x_1^2+\cdots+x_n^2}$. 对 $i=1,\cdots,n$，则
\[
\frac{\partial r}{\partial x_i}=\frac{x_i}{r}.
\]
由链式法则得
\[
\frac{\partial f}{\partial x_i}(x)
= g'(r)\frac{\partial r}{\partial x_i}
= g'(r)\frac{x_i}{r}.
\]
因此
\[
\nabla f(x)=\left(g'(r)\frac{x_1}{r},\cdots,g'(r)\frac{x_n}{r}\right)
= g'(r)\frac{x}{r}.
\]
两边取欧氏范数，得
\[
|\nabla f(x)|=|g'(r)|\left|\frac{x}{r}\right|=|g'(r)|.
\]
再由球坐标变换可知
\begin{align*}
\iiint_{\left| y \right|\geqslant r}{|\nabla f(y)|\,\mathrm{d}y}=\int_r^{+\infty}{\int_0^{2\pi}{\int_0^{\pi}{|g' (s)|s^2\sin \theta \,\mathrm{d}\theta \,\mathrm{d}\phi \,\mathrm{d}s}}}=4\pi \int_0^{+\infty}{|g' (s)|s^2\,\mathrm{d}s}.
\end{align*}
进而
\begin{align}
\int_r^{+\infty}|g'(s)|s^2\,\mathrm{d}s
=\frac{1}{4\pi}\iiint_{|y|\geqslant r}|\nabla f(y)|\,\mathrm{d}y
\leqslant \frac{1}{4\pi}\iiint_{\mathbb{R}^3}|\nabla f(y)|\,\mathrm{d}y,\label{eq::iojifjioJIJIJw-2}
\end{align}
所以由\eqref{eq::iojifjioJIJIJw-1}\eqref{eq::iojifjioJIJIJw-2}式可得
\[
|x|^2|f(x)|=r^2|g(r)|\leqslant \frac{1}{4\pi}\iiint_{\mathbb{R}^3}|\nabla f(y)|\,\mathrm{d}y.
\]
特别地,当 $x=0$ 时，上式依然成立。因此，取
$C=\frac{1}{4\pi},$
则
\[
\max_{x\in\mathbb{R}^3}\bigl(|x|^2|f(x)|\bigr)
\leqslant C\iiint_{\mathbb{R}^3}|\nabla f(x)|\,\mathrm{d}x.
\]\end{proof}

\begin{example}
设 \(\Omega = \{(x,y,z) \in \mathbb{R}^3 : x^2+y^2+z^2 < 1\}\) 且 \(f \in C^2 (\overline{\Omega})\) 满足 \(\Delta f = \sqrt{x^2+y^2+z^2}\), 计算
\begin{align*}
\iiint_{\Omega} \left( x \frac{\partial f}{\partial x} + y \frac{\partial f}{\partial y} + z \frac{\partial f}{\partial z} \right) \mathrm{d}x\mathrm{d}y\mathrm{d}z.
\end{align*}
\end{example}
\begin{solution}
记
\begin{align*}
I = \iiint_{\Omega} \left( x \frac{\partial f}{\partial x} + y \frac{\partial f}{\partial y} + z \frac{\partial f}{\partial z} \right) \mathrm{d}x\mathrm{d}y\mathrm{d}z = \iiint_{\Omega} \nabla f \cdot \vec{r} \, \mathrm{d}V,
\end{align*}
其中 \(\vec{r} = (x,y,z)\) .
应用\hyperref[theorem:Green恒等式]{Green第一恒等式}得
\begin{align*}
\iiint_{\Omega} (u \Delta v + \nabla u \cdot \nabla v) \, \mathrm{d}V = \iint_{\partial \Omega} u \frac{\partial v}{\partial \vec{n}} \, \mathrm{d}S.
\end{align*}
令 \(v = f\), 并选取 \(u = \frac{x^2+y^2+z^2}{2} = \frac{r^2}{2}\),其中$r=\sqrt{x^2+y^2+z^2}$. 则有 \(\nabla u = x\vec{i} + y\vec{j} + z\vec{k} = \vec{r}\)。将其代入上式得
\begin{align*}
\iiint_{\Omega} \left( \frac{r^2}{2} \Delta f + \nabla f \cdot \vec{r} \right) \mathrm{d}V = \iint_{\partial \Omega} \frac{1}{2} \frac{\partial f}{\partial \vec{n}} \mathrm{d}S.
\end{align*}
由此可以推出
\begin{align*}
I = \frac{1}{2} \iint_{\partial \Omega} \frac{\partial f}{\partial \vec{n}} \mathrm{d}S - \frac{1}{2} \iiint_{\Omega} r^2 \Delta f \, \mathrm{d}V.
\end{align*}
接下来分别计算上式中的两项:
\begin{enumerate}[(1)]
\item 对于边界积分, 
由题设知\(\Delta f = r\), 再由\hyperref[theorem:Green恒等式]{Green第一恒等式}和球坐标变换可得
\begin{align*}
\frac{1}{2}\iint_{\partial \Omega}{\frac{\partial f}{\partial \vec{n}}\mathrm{d}S}=\frac{1}{2}\iiint_{\Omega}{\Delta f\,\mathrm{d}V}=\frac{1}{2}\iiint_{\Omega}{r\,\mathrm{d}V}=\frac{1}{2}\int_0^{2\pi}{\mathrm{d}\phi \int_0^{\pi}{\sin \theta \mathrm{d}\theta \int_0^1{r^3\mathrm{d}r}}}=\frac{\pi}{2}.
\end{align*}

\item 对于三重积分,由题设知\(\Delta f = r\),再利用球坐标变换计算得
\begin{align*}
\frac{1}{2} \iiint_{\Omega} r^2 \Delta f \, \mathrm{d}V = \frac{1}{2} \iiint_{\Omega} r^2 \cdot r \, \mathrm{d}V = \frac{1}{2} \int_{0}^{2\pi} \mathrm{d}\phi \int_{0}^{\pi} \sin\theta \mathrm{d}\theta \int_{0}^{1} r^5 \mathrm{d}r = \frac{\pi}{3}.
\end{align*}
\end{enumerate}
将上面两项的结果代回 \(I\) 的表达式, 即可得到
\begin{align*}
I = \frac{\pi}{2} - \frac{\pi}{3} = \frac{\pi}{6}.
\end{align*}
\end{solution}





















\end{document}